\section{Technical Derivations}
\label{app:technical-derivations}

This appendix records notation conventions used in the empirical chapters. The derivations support response-path reporting and do not add structural identifying assumptions.

\subsection{Moving-Average Representation}

Starting from the reduced-form VAR in \Cref{eq:reduced-form-var}, the companion-form representation implies:

\begin{equation}
  \vect{Y}_t = \mat{F}\vect{Y}_{t-1} + \vect{v}_t,
  \label{eq:companion-form}
\end{equation}

where \(\vect{Y}_t\) stacks current and lagged values of \(\vect{y}_t\). The moving-average coefficient at horizon \(h\) is obtained from \(\mat{F}^h\).

\subsection{Cumulative Responses}

For response \(i\), the cumulative response through reported horizon \(H\) is:

\begin{equation}
  \operatorname{CIRF}_i(H) = \sum_{h \in \mathcal{H}: h \leq H}\widehat{\IRF}_i(h),
  \label{eq:cumulative-irf}
\end{equation}

where \(\mathcal{H}=\{0,1,3,6,12,24\}\) is the thesis-facing horizon set. This convention matches the cumulative-response summaries exported by the empirical pipeline and used in the persistence tables.
